% 04_macros.tex
% Componentes de Control de Calidad (Test Cases, Bug Reports, Trazabilidad, Métricas)
% Documento de Pruebas y Validaciones (V&V)

% =====================================================================
% 0. PARCHE BABEL-ESPAÑOL: el guion activo de babel-spanish ejecuta
% \char\hyphenchar\font; en la familia tipográfica monoespaciada (Courier,
% hyphenchar=-1) esto produce "Bad character code (-1)". Se redefine para
% emitir un guion literal seguro (catcode 12) sin recursión.
% =====================================================================
\makeatletter
\@ifundefined{es@chf}{}{\renewcommand{\es@chf}{\string-}}
\makeatother

% =====================================================================
% 1. FORMATO EN LÍNEA: comandos, rutas, endpoints
% =====================================================================
\newcommand{\comando}[1]{{\fontfamily{pcr}\fontsize{9}{10.8}\selectfont\colorbox{grisTecnico}{\texttt{#1}}}}
\newcommand{\ruta}[1]{{\ttfamily\fontsize{8.5}{10.2}\selectfont\seqsplit{#1}}}

\makeatletter
\newcommand{\pathbreak}[1]{\texttt{\@pathbreak#1\relax}}
\def\@pathbreak#1{%
  \ifx#1\relax\else
    \ifx#1/\,\discretionary{}{}{}/\discretionary{}{}{}%
    \else\ifx#1_\_\discretionary{}{}{}%
    \else\ifx#1-#1\discretionary{}{}{}%
    \else\ifx#1.#1\discretionary{}{}{}%
    \else#1\fi\fi\fi\fi
    \expandafter\@pathbreak
  \fi}
\makeatother

% Tipos de columna con centrado vertical
\newcolumntype{L}[1]{>{\raggedright\arraybackslash}m{#1}}
\newcolumntype{C}[1]{>{\centering\arraybackslash}m{#1}}
\newcolumntype{R}[1]{>{\raggedleft\arraybackslash}m{#1}}
\newcolumntype{Y}{>{\raggedright\arraybackslash}X}
\newcolumntype{Z}{>{\centering\arraybackslash}X}

% =====================================================================
% 2. BADGES DE ESTADO Y SEVERIDAD
% =====================================================================
\newcommand{\estadoPass}{\cellcolor{bgPass}{\sffamily\bfseries\color{white}PASS}}
\newcommand{\estadoFail}{\cellcolor{bgFail}{\sffamily\bfseries\color{white}FAIL}}
\newcommand{\estadoBlocked}{\cellcolor{bgBlocked}{\sffamily\bfseries\color{white}BLOCK}}
\newcommand{\estadoSkip}{\cellcolor{bgSkip}{\sffamily\bfseries\color{white}SKIP}}

% Badges en línea (para texto corrido y celdas)
\newcommand{\badge}[2]{{\setlength{\fboxsep}{2pt}\colorbox{#1}{\sffamily\bfseries\fontsize{7.5}{9}\selectfont\color{white}\,#2\,}}}
\newcommand{\sevCrit}{\badge{sevCritica}{CRÍTICA}}
\newcommand{\sevAlt}{\badge{sevAlta}{ALTA}}
\newcommand{\sevMed}{\badge{sevMedia}{MEDIA}}
\newcommand{\sevBaj}{\badge{sevBaja}{BAJA}}
\newcommand{\bPass}{\badge{bgPass}{PASS}}
\newcommand{\bFail}{\badge{bgFail}{FAIL}}
\newcommand{\bBlock}{\badge{bgBlocked}{BLOCK}}
% Estado de resolución de defectos
\newcommand{\resCerrado}{\badge{bgPass}{CERRADO}}
\newcommand{\resAbierto}{\badge{sevAlta}{ABIERTO}}
\newcommand{\resResuelto}{\badge{sevBaja}{RESUELTO}}
\newcommand{\resVerif}{\badge{bgPass}{VERIFICADO}}
\newcommand{\resDiferido}{\badge{bgSkip}{DIFERIDO}}

% =====================================================================
% 3. TABLA-CATÁLOGO DE CASOS DE PRUEBA (resumen de un sprint, longtable)
% Uso: \begin{CatalogoTC} ... \tcRow{ID}{Título}{Tipo}{Prio}{\bPass} ... \end{CatalogoTC}
% =====================================================================
\newenvironment{CatalogoTC}{%
    \renewcommand{\arraystretch}{1.3}%
    \fontsize{8}{9.6}\selectfont\sffamily%
    \setlength{\tabcolsep}{4pt}%
    \begin{longtable}{%
        |>{\ttfamily\bfseries\color{colorPrimario}}p{1.9cm}%
        |>{\raggedright\arraybackslash}p{6.4cm}%
        |>{\centering\arraybackslash}p{2.4cm}%
        |>{\centering\arraybackslash}p{1.0cm}%
        |>{\centering\arraybackslash}p{1.4cm}|}
    \hline
    \rowcolor{headerTabla}
    \sffamily\textbf{ID} & \sffamily\textbf{Título del caso} & \sffamily\textbf{Tipo} & \sffamily\textbf{Prio} & \sffamily\textbf{Result.} \\
    \hline
    \endfirsthead
    \multicolumn{5}{l}{\sffamily\footnotesize\itshape\color{grisISO}Catálogo de casos de prueba (continuación)}\\
    \hline
    \rowcolor{headerTabla}
    \sffamily\textbf{ID} & \sffamily\textbf{Título del caso} & \sffamily\textbf{Tipo} & \sffamily\textbf{Prio} & \sffamily\textbf{Result.} \\
    \hline
    \endhead
    \hline
    \multicolumn{5}{r}{\sffamily\tiny\color{grisISO}\textit{continúa en la página siguiente}}\\
    \endfoot
    \hline
    \endlastfoot
}{%
    \end{longtable}%
}
\newcommand{\tcRow}[5]{#1 & #2 & #3 & #4 & #5 \\ \hline}

% =====================================================================
% 4. CASO DE PRUEBA DETALLADO (ficha forense)
% Diseño: barra de título de color + tablas tabularx INDEPENDIENTES en el
% flujo del texto. No se anida tabularx dentro de un tcolorbox breakable
% (eso provoca cascadas de error de "TX@get@body"); en cambio, cada bloque es
% una tabla pequeña separada, por lo que LaTeX puede quebrar entre bloques.
% Uso:
% \begin{FichaTC}{ID}{Título} ... \end{FichaTC}
% =====================================================================
\newcommand{\fichaTitulo}[3]{% color, icono, texto
    \vspace{0.35cm}\par\noindent
    {\setlength{\fboxsep}{4pt}\noindent\colorbox{#1}{\parbox{\dimexpr\linewidth-8pt\relax}{%
       \sffamily\bfseries\small\color{white}#2\ \ #3}}}\par\nopagebreak\vspace{1pt}}

\newenvironment{FichaTC}[2]{%
    \fichaTitulo{colorPrimario}{\faVial}{#1 — #2}%
    \renewcommand{\arraystretch}{1.2}\fontsize{8.5}{10.4}\selectfont\sffamily
}{%
    \vspace{0.2cm}}
\newcommand{\tcMeta}[4]{%
    \noindent\begin{tabularx}{\linewidth}{|>{\bfseries\color{colorPrimario}}p{2.4cm}|Y|>{\bfseries\color{colorPrimario}}p{2.0cm}|Y|}
    \hline\rowcolor{grisTecnico}
    Tipo & #1 & Prioridad & #2 \\ \hline
    \rowcolor{grisTecnico}
    Componente & #3 & Requisito & #4 \\ \hline
    \end{tabularx}\par}
\newcommand{\tcCampo}[2]{%
    \noindent\begin{tabularx}{\linewidth}{|>{\bfseries\color{colorPrimario}}p{3.2cm}|Y|}
    \hline #1 & #2 \\ \hline \end{tabularx}\par}
% Tabla de pasos enumerados
\newenvironment{tcPasos}{%
    \noindent\tabularx{\linewidth}{|c|Y|Y|}
    \hline\rowcolor{headerTabla}
    \sffamily\textbf{\#} & \sffamily\textbf{Acción (instrucción)} & \sffamily\textbf{Resultado esperado} \\ \hline
}{%
    \endtabularx\par}
\newcounter{tcpasocnt}
\newcommand{\paso}[2]{\stepcounter{tcpasocnt}\arabic{tcpasocnt} & #1 & #2 \\ \hline}
\newcommand{\resetPasos}{\setcounter{tcpasocnt}{0}}
\newcommand{\tcResultado}[3]{%
    \noindent\begin{tabularx}{\linewidth}{|>{\bfseries\color{colorPrimario}}p{3.2cm}|Y|}
    \hline\rowcolor{headerTabla}
    \multicolumn{2}{|l|}{\sffamily\textbf{Veredicto de ejecución}} \\ \hline
    Resultado esperado & #1 \\ \hline
    Resultado real & #2 \\ \hline
    Estado & #3 \\ \hline
    \end{tabularx}\par}

% =====================================================================
% 5. REPORTE DE DEFECTO (BUG REPORT) FORENSE
% =====================================================================
\newenvironment{FichaBug}[2]{%
    \fichaTitulo{bgFail}{\faBug}{#1 — #2}%
    \renewcommand{\arraystretch}{1.2}\fontsize{8.5}{10.4}\selectfont\sffamily
}{%
    \vspace{0.2cm}}
\newcommand{\bugMeta}[4]{%
    \noindent\begin{tabularx}{\linewidth}{|>{\bfseries\color{bgFail}}p{2.2cm}|Y|>{\bfseries\color{bgFail}}p{1.8cm}|Y|}
    \hline\rowcolor{grisTecnico}
    Severidad & #1 & Prioridad & #2 \\ \hline
    \rowcolor{grisTecnico}
    Componente & #3 & Estado & #4 \\ \hline
    \end{tabularx}\par}
\newcommand{\bugCampo}[2]{%
    \noindent\begin{tabularx}{\linewidth}{|>{\bfseries\color{bgFail}}p{3.4cm}|Y|}
    \hline #1 & #2 \\ \hline \end{tabularx}\par}
% Bloque ancho (para pasos, RCA, evidencia) con título
\newcommand{\bugBloque}[2]{%
    \noindent\begin{tabularx}{\linewidth}{|Y|}
    \hline\rowcolor{headerTabla}\sffamily\textbf{#1} \\ \hline
    #2 \\ \hline \end{tabularx}\par}

% =====================================================================
% 6. TABLA DE PRUEBAS DE API (endpoint / payload / código HTTP)
% =====================================================================
\newenvironment{TablaAPI}{%
    \renewcommand{\arraystretch}{1.35}\vspace{0.25cm}\noindent
    \fontsize{8.5}{10.2}\selectfont\sffamily
    \tabularx{\textwidth}{|>{\bfseries\color{colorPrimario}}p{3.2cm}|Y|}
    \hline\rowcolor{headerTabla}
    \sffamily\textbf{Atributo de la petición} & \sffamily\textbf{Valor} \\ \hline
}{%
    \endtabularx\vspace{0.25cm}}

% =====================================================================
% 7. MATRIZ DE TRAZABILIDAD (RTM) — longtable
% =====================================================================
\newenvironment{MatrizRTM}{%
    \renewcommand{\arraystretch}{1.3}%
    \fontsize{8}{9.6}\selectfont\sffamily\setlength{\tabcolsep}{4pt}%
    \begin{longtable}{%
        |>{\ttfamily\bfseries\color{colorPrimario}}p{2.0cm}%
        |>{\raggedright\arraybackslash}p{5.3cm}%
        |>{\ttfamily\raggedright\arraybackslash}p{3.0cm}%
        |>{\ttfamily\raggedright\arraybackslash}p{2.6cm}|}
    \hline\rowcolor{headerTabla}
    \sffamily\textbf{Requisito} & \sffamily\textbf{Historia / Descripción} & \sffamily\textbf{Casos de prueba} & \sffamily\textbf{Defectos} \\ \hline
    \endfirsthead
    \multicolumn{4}{l}{\sffamily\footnotesize\itshape\color{grisISO}Matriz de trazabilidad (continuación)}\\ \hline
    \rowcolor{headerTabla}
    \sffamily\textbf{Requisito} & \sffamily\textbf{Historia / Descripción} & \sffamily\textbf{Casos de prueba} & \sffamily\textbf{Defectos} \\ \hline
    \endhead
    \hline\multicolumn{4}{r}{\sffamily\tiny\color{grisISO}\textit{continúa en la página siguiente}}\\ \endfoot
    \hline\endlastfoot
}{%
    \end{longtable}}
\newcommand{\rtmRow}[4]{#1 & #2 & #3 & #4 \\ \hline}

% =====================================================================
% 8. GRÁFICO DE BARRAS DE EJECUCIÓN (TikZ) — pasaron vs fallaron
% \graficoEjecucion{pasaron}{fallaron}{bloqueados}{total}
% =====================================================================
\newcommand{\graficoEjecucion}[4]{%
\begin{center}
\begin{tikzpicture}[font=\sffamily\scriptsize]
  \def\maxv{#4}
  \pgfmathsetmacro{\sc}{6.0/\maxv}
  \pgfmathsetmacro{\bp}{#1*\sc}
  \pgfmathsetmacro{\bf}{#2*\sc}
  \pgfmathsetmacro{\bb}{#3*\sc}
  \fill[bgPass] (0,2.0) rectangle (\bp,2.6); \node[right] at (\bp,2.3) {Pasaron: #1};
  \fill[bgFail] (0,1.1) rectangle (\bf,1.7); \node[right] at (\bf,1.4) {Fallaron: #2};
  \fill[bgBlocked] (0,0.2) rectangle (\bb,0.8); \node[right] at (\bb,0.5) {Bloqueados: #3};
  \draw[grisISO,->] (0,0) -- (6.6,0);
  \draw[grisISO,->] (0,0) -- (0,3.0);
  \node[left, text=grisISO] at (0,2.3) {};
\end{tikzpicture}
\end{center}}

% Barra de cobertura horizontal sencilla. \barraCobertura{porcentaje}{etiqueta}
\newcommand{\barraCobertura}[2]{%
\noindent{\sffamily\fontsize{9}{11}\selectfont #2\par}%
\noindent\begin{tikzpicture}[font=\sffamily\scriptsize]
  \pgfmathsetmacro{\w}{#1*0.12}
  \fill[grisTecnico] (0,0) rectangle (12,0.45);
  \fill[colorPrimario] (0,0) rectangle (\w,0.45);
  \draw[grisISO] (0,0) rectangle (12,0.45);
  \node[right, text=white, font=\sffamily\bfseries\scriptsize] at (0.15,0.225) {#1\%};
\end{tikzpicture}\par\vspace{0.2cm}}

% =====================================================================
% 9. CAJAS DE AVISO (reutilizadas del estándar corporativo)
% =====================================================================
\newtcolorbox{NotaTecnica}{
    colback=white, colframe=bordeNotaTecnica, boxrule=1.2pt, arc=2pt,
    title={\faInfoCircle\ \textbf{Nota Técnica}}, coltitle=bordeNotaTecnica,
    fonttitle=\sffamily, attach title to upper=\par\vspace{0.2em}}
\newtcolorbox{AvisoNota}{
    colback=bgNota, colframe=bgNota!80!black, boxrule=0.5pt, leftrule=4pt, arc=2pt,
    title={\faInfoCircle\ \textbf{Nota}}, coltitle=black, fonttitle=\sffamily,
    attach title to upper=\quad}
\newtcolorbox{AlertaDefecto}{
    colback=bgAdvertencia, colframe=red!80!black, boxrule=0.5pt, leftrule=4pt, arc=2pt,
    title={\faBug\ \textbf{Hallazgo Crítico}}, coltitle=black, fonttitle=\sffamily,
    attach title to upper=\par\vspace{0.2em}}
\newtcolorbox{AlertaSeguridad}{
    colback=bgAdvertencia, colframe=red!80!black, boxrule=0.5pt, leftrule=4pt, arc=2pt,
    title={\faShield*\ \textbf{Alerta de Seguridad}}, coltitle=black, fonttitle=\sffamily,
    attach title to upper=\par\vspace{0.2em}}
\newtcolorbox{Veredicto}{
    colback=white, colframe=bgPass, boxrule=1.2pt, arc=2pt,
    title={\faCheckCircle\ \textbf{Veredicto de Aceptación}}, coltitle=bgPass,
    fonttitle=\sffamily, attach title to upper=\par\vspace{0.2em}}

% =====================================================================
% 10. LISTAS COMPACTAS Y CÓDIGO
% =====================================================================
\setlist[itemize]{noitemsep, topsep=2pt}
\setlist[enumerate]{noitemsep, topsep=2pt}
\newlist{checklist}{itemize}{1}
\setlist[checklist]{label=$\square$, leftmargin=*, noitemsep}

\lstdefinelanguage{JavaScript}{
    keywords={typeof, new, true, false, catch, function, return, null, switch, var,
        let, const, if, in, of, while, do, else, case, break, import, export, from,
        await, async, class, extends, default, try, finally, throw, this, void},
    keywordstyle=\color{codeKeyword}\bfseries,
    ndkeywords={Promise, window, sessionStorage, localStorage, JSON, Math, Date, axios},
    ndkeywordstyle=\color{c4Sistema}\bfseries,
    comment=[l]{//}, morecomment=[s]{/*}{*/}, commentstyle=\color{codeComment}\itshape,
    stringstyle=\color{codeString}, morestring=[b]', morestring=[b]", morestring=[b]`,
    sensitive=true}
\lstalias{TypeScript}{JavaScript}
\lstdefinestyle{ethosCode}{
    basicstyle=\ttfamily\fontsize{7.6}{9.2}\selectfont, extendedchars=true,
    backgroundcolor=\color{grisTecnico}, keywordstyle=\color{codeKeyword}\bfseries,
    commentstyle=\color{codeComment}\itshape, stringstyle=\color{codeString},
    numbers=left, numberstyle=\tiny\color{grisISO}, numbersep=8pt,
    frame=single, framesep=5pt, rulecolor=\color{grisISO!50},
    breaklines=true, breakatwhitespace=true, showstringspaces=false, tabsize=2,
    captionpos=b, aboveskip=0.6em, belowskip=0.6em,
    literate=%
        {á}{{\'a}}1 {é}{{\'e}}1 {í}{{\'i}}1 {ó}{{\'o}}1 {ú}{{\'u}}1
        {Á}{{\'A}}1 {É}{{\'E}}1 {Í}{{\'I}}1 {Ó}{{\'O}}1 {Ú}{{\'U}}1
        {ñ}{{\~n}}1 {Ñ}{{\~N}}1 {¿}{{?`}}1 {¡}{{!`}}1}
\lstset{style=ethosCode}
